\documentclass[11pt,letterpaper]{article}
\usepackage[T1]{fontenc}
\usepackage{lmodern}
\usepackage[margin=1.04in]{geometry}
\usepackage{microtype}
\usepackage{amsmath,amssymb,amsthm}
\usepackage{booktabs}
\usepackage{xcolor}
\definecolor{linkblue}{RGB}{30,69,101}
\usepackage[colorlinks=true,linkcolor=linkblue,citecolor=linkblue,urlcolor=linkblue]{hyperref}
\newcommand*{\manuscriptauthor}{Robert Thomas Collins}
\hypersetup{pdftitle={Reducing simple stochastic reachability games to deterministic mean-payoff games},
  pdfauthor={\manuscriptauthor}}

\newtheorem{theorem}{Theorem}
\newtheorem{lemma}[theorem]{Lemma}
\newtheorem{proposition}[theorem]{Proposition}
\newcommand{\Win}{\mathsf{Win}}
\newcommand{\Lose}{\mathsf{Lose}}
\newcommand{\val}{\operatorname{val}}
\newcommand{\bits}{\operatorname{bits}}
\newcommand{\sys}{\mathcal S}
\newcommand{\compiled}{\mathcal C}
\setlength{\emergencystretch}{1.5em}

\begin{document}
\begin{center}
  {\LARGE\bfseries Reducing simple stochastic reachability games\\[.10em]
    to deterministic mean-payoff games\par}
  \medskip{\large\manuscriptauthor\par}
  \medskip{\small Working draft \enspace\textperiodcentered\enspace 24 November 2026\par}
\end{center}

\begin{abstract}
Given a finite simple stochastic reachability game, a start state, and a binary
rational threshold \(t\), we construct one deterministic mean-payoff game with
integer edge weights in polynomial bit time. Its value at the designated start is
nonnegative exactly when the original winning probability under optimal play is at
least \(t\), even when play can continue forever. An explicit error bound lets us
replace fair-coin averages by min/max circuits, including averages that depend on
one another around a cycle.
\end{abstract}

\section{The games and the threshold question}

Max and Min move a token on a finite directed graph. At a state owned by one of
the players, that player chooses the next state from a nonempty list. At a chance
state, a fair coin chooses between two listed successors, which may coincide.
Both players see the moves so far and may use the whole history. The graph
has two terminal states, \(\Win\) and \(\Lose\). The token stays at a terminal
once it gets there. Max earns one if the token ever reaches \(\Win\).
All other plays earn zero, including those that go on forever. This is Condon's simple stochastic reachability game~\cite{condon}.
Given a start \(s\) and a rational threshold \(t\), can Max guarantee a
winning probability of at least \(t\), however Min plays? Equality counts as yes.

In the game we construct, Max or Min chooses the next edge at every vertex.
Every vertex has an outgoing edge. Each move pays Max the integer weight of the
edge traversed. Max wants a high average reward per move and Min wants a low
one. For an infinite play with edges \(e_1,e_2,\ldots\) of weights \(w(e_i)\),
the payoff is the lower limiting average
\[
       \liminf_{k\to\infty}\frac{1}{k}\sum_{i=1}^{k} w(e_i).
\]
The lower limit gives a payoff even when the running averages do not converge.

In either game the value \(\val\) at a given start is the best payoff Max can
secure against Min. In the target game we ask whether Max can secure a
nonnegative payoff. To construct this one game we use a number of bit operations
polynomial in the number of bits it takes to list the source game, its start,
and the numerator and denominator of \(t\).

\begin{theorem}\label{thm:main}
There is a polynomial-bit-time many-one reduction that takes an explicitly given
simple stochastic reachability game \(G\), a state \(s\), and a rational \(t=a/b\)
in binary, with \(b>0\), to a finite deterministic mean-payoff game \(H\) with
binary integer edge weights and a state \(s_0\), such that
\[
                 \val_G(s)\ge t \quad\Longleftrightarrow\quad \val_H(s_0)\ge 0 .
\]
\end{theorem}

Write \(P(\sigma,\tau)\) for the payoff at a given start when Max uses strategy
\(\sigma\) and Min uses \(\tau\), taking the win probability in the stochastic game.
The lower value is \(\sup_\sigma\inf_\tau P(\sigma,\tau)\) and the upper value is
\(\inf_\tau\sup_\sigma P(\sigma,\tau)\). Their common value is \(\val\).
Finite deterministic mean-payoff games admit optimal \emph{positional} strategies,
which fix one outgoing edge at each vertex a player owns~\cite{em,gkk}.

Together with the reduction from deterministic mean-payoff games to simple
stochastic games by Zwick and Paterson~\cite[\S\S5--6]{zp}, Theorem~\ref{thm:main}
makes the two threshold decision problems polynomial-time equivalent. The abstract
of Parys's Warsaw seminar of 28 October 2026 records the direction from simple
stochastic games to mean-payoff games as open at that date~\cite{parys}.

We begin by asking how Min could certify that \(\val_G(s)<t\). Fix the
terminal labels at one for \(\Win\) and zero for \(\Lose\). Give each other
state a label \(x_i\) between zero and one. Take the largest successor label
at a Max state, the smallest at a Min state, and the average at a chance
state. If every result is at most the current label, Lemma~\ref{lem:reach}
shows that each \(x_i\) bounds its state's winning probability from above. Then \(x_s<t\) certifies
a no answer.
Now replace each label by its complement \(1-x_i\). At a chance state, the
average of the two complementary successor labels is an upper bound on the
current label. Replacing the average by a larger number loosens the bound.
If two chance states are successors of each other, raising either label can
loosen the bound on the other. The example after Proposition~\ref{prop:compile}
shows how a margin that is too small lets a false certificate pass at a
threshold of \(1/3\).

The construction uses two bounds computed from the input. First we fix a
positive margin that a certificate can meet whenever the value is strictly
below \(t\). Next, fix a supporting move at each Min state. If those choices
can pass only by overestimating some average, how small could the largest
overestimate be? We give a positive lower bound that works for every such
choice. We keep each circuit's error below it, so a certificate can pass
only if exact labels exist for the chosen moves. An intermediate game \(K\)
tests for such a certificate. Changing its owners and weights then reverses
its answer.

Bodirsky, Loho, and Skomra describe inequalities using maxima and averages for
the \(1/2\)-threshold in a stopping stochastic game, one in which absorption
occurs almost surely under every pair of strategies~\cite[Prop.~5]{bls}.
Atserias and Maneva relate inequalities using maxima and integer offsets to
mean-payoff games~\cite{am}. To pass between these two kinds of inequalities
we use the circuit layers of Bezem, Nieuwenhuis, and Rodr\'iguez-Carbonell and
their one-sided inequality formulation by Allamigeon, Benchimol, Gaubert, and
Joswig~\cite[\S3]{bnr}\cite[Prop.~20]{abgj}. Our precision bound determines how
many layers suffice for all the averages at once.

\section{Certifying a reachability value below the threshold}
\label{sec:reach}

Write \(n\) for the number of nonterminal states. If \(t\le 0\), the answer is yes.
If \(t>1\), the answer is no. For a terminal start, compare its known value with \(t\).
In each case, output a single vertex with a loop of weight zero for yes or \(-1\) for
no. Henceforth \(s\) is nonterminal and \(0<t\le 1\).

For \(x\in[0,1]^n\), extend \(x\) to the terminals by \(x_{\Win}=1\) and \(x_{\Lose}=0\).
The Bellman operator \(F\) takes the maximum of the successor labels at a Max state,
the minimum at a Min state, and the arithmetic mean at a chance state. All vector
inequalities below are coordinatewise. If a nonterminal state's only move
leads back to itself, every label satisfies its Bellman equation. Its actual
winning probability is zero. The least fixed point singles out the reachability
value.

\begin{lemma}\label{lem:reach}
The reachability value vector \(v\) is the least fixed point of \(F\) on the unit cube.
In particular,
\[
 v_s<t
 \quad\Longleftrightarrow\quad
 \exists\,x\in[0,1]^n\quad F(x)\le x,\quad x_s<t .
\]
\end{lemma}

\begin{proof}
Let \(u_0=0\) on nonterminals and \(u_{k+1}=F(u_k)\), always supplying the fixed
terminal labels. Backward induction identifies \(u_k(i)\) with the value for reaching
\(\Win\) within \(k\) moves from \(i\). A finite-horizon Max strategy guarantees
this much probability of eventual reachability as well. The vectors \(u_k\) increase
to a limit \(u\) in the cube. Finite maxima, minima, and means are continuous, so
\(F(u)=u\).

Fix one vector \(x\) with \(F(x)\le x\). All finite-horizon values stay below it,
since \(u_0\le x\) and monotonicity gives
\[
                     u_{k+1}=F(u_k)\le F(x)\le x .
\]
There is also a single Min strategy giving the upper bound \(x_i\) against every
Max strategy for eventual reachability. At each Min state choose a successor of
minimum \(x\)-label. Then the conditional expected next label is at most the present
label at every state. For the state \(I_k\) after \(k\) moves from \(i\), absorption at
\(\Win\) and nonnegativity elsewhere give
\[
    \Pr_i(\Win\text{ reached by time }k)
    \ \le\ \mathbb E_i[x_{I_k}]
    \ \le\ x_i .
\]
The events on the left increase to eventual reachability. Taking \(x=u\) bounds the
upper value by \(u_i\), while the finite-horizon Max guarantees give a lower
value of at least \(u_k(i)\) for every \(k\). Thus \(v=u\), and every \(x\) with \(F(x)\le x\)
dominates \(v\). When \(v_s<t\), \(x=v\) supplies the required witness.
\end{proof}

Set \(y_i=1-x_i\), with terminal literals \(y_{\Win}=0\) and \(y_{\Lose}=1\).
The \(x_i\) certify upper bounds on winning probabilities. The corresponding \(y_i\)
give lower bounds on their complements, which include the possibility of never
absorbing. In these coordinates the certificate consists of \(0\le y_i\le1\),
\(y_s>1-t\), and the following rows. A row is one inequality in a simultaneous
system.
\[
\begin{array}{@{}ll@{}}
\toprule
\text{State \(i\) in \(G\)} & \text{Constraint on the labels}\\
\midrule
\text{Max, with successors \(J_i\)} & y_i\le y_j\quad\text{for every }j\in J_i\\[2pt]
\text{Min, with successors \(J_i\)} & y_i\le\max_{j\in J_i} y_j\\[2pt]
\text{Chance, with successors \(j,k\)} & y_i\le(y_j+y_k)/2\\
\bottomrule
\end{array}
\]
At a source Max state an upper bound must cover every choice Max might make. At a
Min state one successor supporting the bound suffices, because Min can choose it.

\section{How much precision do the inequalities need?}
\label{sec:precision}

A \emph{branch} chooses one term \(y_j\) in every right-hand maximum and imposes
\(y_i\le y_j\) in its place. Every solution satisfies a branch, by choosing terms
attaining the maxima. Every solution of any branch satisfies the original rows.
Thus we can reason about linear inequalities branch by branch. The following
estimate will bound both the margin at the start state and the error in the
chance rows. The input bounds the row norms and denominators for all branches
at once.

\begin{lemma}\label{lem:vertex}
Let a nonempty bounded polyhedron in \(\mathbb R^d\) be given by weak inequalities
whose coefficient rows are integral with Euclidean norm at most \(A\). Suppose a
positive integer \(B\) clears all right-hand sides. If the minimum or maximum of a
coordinate is positive, it is at least \(1/(B A^{d})\).
\end{lemma}

\begin{proof}
The optimal face contains a vertex. Its tight rows span \(\mathbb R^d\), since
otherwise a small segment in a common null direction would stay feasible on both
sides of the vertex. Choose \(d\) independent tight rows and write their equalities
as \(Mz=q\). The \emph{coefficients} of \(M\) are integers, while \(q\) and the
coordinates of the vertex may be fractional. Hadamard's inequality gives
\(0<|\det M|\le A^d\). Since \(Bq\) is integral, Cramer's rule gives
\(B\det(M)z_j\in\mathbb Z\) for each \(j\). Apply this to the positive optimal
coordinate. The argument also applies to a polyhedron of smaller affine dimension.
\end{proof}

\subsection{A margin for the strict threshold}

A strict witness puts the start label above \(1-t\). We can replace the strict
test by a weak inequality if we know how far above \(1-t\) at least one
witness can reach. Put
\[
              D=b\,3^{\,n+1},\qquad \delta=\frac1D,
              \qquad \theta=1-t+\delta .
\]
Let \(\sys_\delta\) be the system in Section~\ref{sec:reach} with the single strict
row replaced by \(y_s\ge\theta\). We claim that
\begin{equation}\label{eq:tighten}
 v_s<t \quad\Longleftrightarrow\quad \sys_\delta\text{ has a real solution}
       \quad\Longleftrightarrow\quad \sys_\delta\text{ has a rational solution}.
\end{equation}
A strict real witness supplies a branch. On this branch retain the box and state
rows, adjoin \(0\le\lambda\le1\), and maximize \(\lambda\) under
\(y_s\ge1-t+\lambda\). The witness gives a feasible \(\lambda>0\).
This polytope is bounded in the \(n+1\)
coordinates \((y,\lambda)\). Clear each average row by multiplying by two, moving
the terminal literals to the right. Its coefficient row has norm at most
\(\sqrt8\), since the left coefficient is \(2\) and a repeated successor can have
coefficient \(-2\). Combining an occurrence of the left variable on the right
only decreases this bound. All other rows have norm at most \(\sqrt2\), and
all right-hand denominators divide \(b\). Lemma~\ref{lem:vertex} with \(A=3\)
shows that the maximum is at least \(1/D\). At the optimal vertex used in that
lemma, Cramer's rule makes \emph{every} coordinate rational. Its source labels
therefore give a rational witness on this branch that survives the tightening.
Every rational witness is real, and every real solution of \(\sys_\delta\) has
\(y_s>1-t\) because \(\delta>0\). This proves~\eqref{eq:tighten}.

\subsection{The error needed for an infeasible branch to pass}

Fix any branch of \(\sys_\delta\). Keep the threshold, box, and all player rows
exact. Give every chance row the same slack \(\rho\) and minimize \(0\le\rho\le1\)
subject to the resulting inequalities
\[
                        y_i\le (y_j+y_k)/2+\rho .
\]
If the relaxed branch is feasible, its minimum \(\rho_{\min}\) is attained on a
compact polytope. The right-hand denominators now divide \(D\), since \(b\)
divides \(D\). In the \(n+1\) coordinates \((y,\rho)\), the doubled average row
can have coefficients \(2,-2,-2\) on the left variable, a repeated successor, and
\(\rho\). Its norm is at most \(\sqrt{12}<4\), including when some of the variables
coincide or an operand is a terminal. Lemma~\ref{lem:vertex} gives the uniform
alternative
\begin{equation}\label{eq:gap}
       \rho_{\min}=0\qquad\text{or}\qquad
       \rho_{\min}\ge\gamma,\qquad
       \gamma=\frac{1}{D\,4^{\,n+1}} .
\end{equation}
With \(\bits(D)=1+\lfloor\log_2D\rfloor\), choose
\begin{equation}\label{eq:depth}
       r=\bits(D)+2(n+1),\qquad \varepsilon=2^{-r-1}<\gamma .
\end{equation}
The margin \(\delta\) stays fixed while we allow an error of at most
\(\varepsilon\) in each average. Since \(\varepsilon<\gamma\), any branch that
passes these relaxed tests attains zero slack.

\section{Eliminating the averages}

We can now replace the average in each chance row by a circuit of depth
determined by~\eqref{eq:depth}. The circuit keeps the average between two
outputs whose separation we can bound after each round. We use the upper
output in the inequality.

The layers below are a recentered form of a min/max construction of Bezem,
Nieuwenhuis, and Rodr\'iguez-Carbonell~\cite[\S3]{bnr}. In the formulation of
Allamigeon, Benchimol, Gaubert, and Joswig, subtracting the common stage translation
\(1-2^{-j-1}\) from the two outputs gives the shifts below. Those authors also
use upper inequalities and monotonicity to recover the greatest gate values at fixed
inputs~\cite[Prop.~20 and eq.~(30)]{abgj}.

For \(u,v\in[0,1]\), put \(U_0=\max(u,v)\), \(L_0=\min(u,v)\). At round
\(j=1,\ldots,r\), set
\[
 \begin{aligned}
  U_j&=\max\{U_{j-1}-2^{-j-1},\,L_{j-1}+2^{-j-1}\},\\
  L_j&=\min\{U_{j-1}-2^{-j-1},\,L_{j-1}+2^{-j-1}\}.
 \end{aligned}
\]
The sum of the two outputs stays \(u+v\). Their gap \(d_j=U_j-L_j\) obeys
\(d_j=|d_{j-1}-2^{-j}|\). Starting with \(d_0\le1\), the bound
\(0\le d_{j-1}\le2\cdot 2^{-j}\) implies \(d_j\le2^{-j}\). The upper bound
halves at every round even when the actual gap increases, as it does in the
first round for two equal inputs. Writing \(m=(u+v)/2\) yields
\begin{equation}\label{eq:circuit}
                 m\le U_r(u,v)=m+d_r/2\le m+\varepsilon .
\end{equation}
Both gates at a round read the same two earlier outputs. With shared wires the
circuit has \(2+2r\) gates.

Replace each average row \(y_i\le(y_j+y_k)/2\) of \(\sys_\delta\) by a fresh
copy of this circuit on \(y_j,y_k\). For each gate make a variable \(z\).
If its two translated input terms are \(p,q\), encode an upper gate with the
single row \(z\le\max(p,q)\), and a lower gate with the two rows \(z\le p\)
and \(z\le q\). Impose \(y_i\le z_{U_r}\), where \(z_{U_r}\) is the final upper
gate variable in this copy. Denote the resulting system by \(\compiled_r\).

The compiled rows allow a source label to exceed the exact average at its state.
To recover an exact witness, we may have to change the source labels together.
The common error bound lets us choose new labels on the same branch.

\begin{proposition}\label{prop:compile}
The tightened exact system \(\sys_\delta\) and the compiled system
\(\compiled_r\) are feasible in exactly the same instances.
\end{proposition}

\begin{proof}
\emph{Every exact witness extends with its source labels unchanged.}
Let \(y\) satisfy \(\sys_\delta\). Evaluate every gate exactly at the inputs
given by \(y\), and assign the result to its gate variable. This satisfies all
upper and lower gate rows. For each replaced row, exact feasibility and
the left side of~\eqref{eq:circuit} give
\[
                y_i\ \le\ \frac{y_j+y_k}{2}\ \le\ U_r(y_j,y_k)=z_{U_r}.
\]
Thus this same \(y\), together with its evaluated gates, satisfies \(\compiled_r\).

\emph{Every compiled witness selects a branch with an exact witness.}
Let \((\hat y,\hat z)\) satisfy \(\compiled_r\), and fix its source labels \(\hat y\).
In each circuit the gate rows bound every assigned gate value above by its exact
evaluation on these inputs. Indeed, this holds for the first gates, and each translated
maximum or minimum is monotone on all real operands. Induction in circuit order
therefore gives \(\hat z_{U_r}\le U_r(\hat y_j,\hat y_k)\), even if some auxiliary
labels lie outside \([0,1]\). Dependencies between source states may be cyclic,
since the induction evaluates only the gates within a single copy at fixed \(\hat y\).
It follows from~\eqref{eq:circuit} that every original average row holds at
\(\hat y\) with slack \(\varepsilon\).

At \(\hat y\), choose an attaining term in each \emph{original} right-hand maximum.
All selected player rows, box bounds, and the tightened threshold hold exactly.
The resulting branch thus admits \((\hat y,\varepsilon)\), so
\(0\le\rho_{\min}\le\varepsilon<\gamma\). By~\eqref{eq:gap} its minimum is
zero, and it is attained. An optimal vertex \((y^*,0)\) supplies exact rational
source labels on this branch. Its coordinates \(y^*\) may differ from \(\hat y\).
\end{proof}

\subsection{Equality on a nonterminating example}

Let \(S\) be a fair coin with successors \(T,Z\), and \(T\) a fair coin with
successors \(S,\Win\). Let \(Z\) have only a move to itself. A two-step excursion
from \(S\) reaches \(\Win\) with probability \(1/4\) and returns to \(S\) with
probability \(1/4\). Thus \(v_S=1/4+v_S/4\), and
\((v_S,v_T,v_Z)=(1/3,2/3,0)\). The vector \((1,1,1)\) is also a Bellman
fixed point, since the equality at \(Z\) allows any label there.

At \(s=S,\ t=1/3\), we have \(D=243\) and require
\(y_S\ge 2/3+1/243=163/243\). The relaxed chance rows interact through the
return edge. Using \(y_Z\le1\) and \(y_{\Win}=0\) gives
\[
 y_T\le y_S/2+\rho,\qquad
 y_S\le (y_T+1)/2+\rho
      \le y_S/4+1/2+3\rho/2 .
\]
Consequently \(y_S\le2/3+2\rho\). Reaching the imposed threshold requires
\(\rho\ge1/486\), and \((y_S,y_T,y_Z)=(163/243,82/243,1)\) attains it.
The input bound gives \(\gamma=1/62{,}208\). The chosen \(r=16\) gives
\(\varepsilon=1/131{,}072\), too small for the two approximate rows together
to create a witness. The strict certificate fails, as equality in the source
requires.

At the same depth, inputs \(1,1/3\) give \(d_j=2^{-j}/3\) for \(j\ge1\) by the
gap recurrence, so \(U_{16}(1,1/3)=2/3+1/393{,}216\). If the margin is reduced
to \(\delta'=1/786{,}432\), the labels \((2/3+\delta',1/3,1)\) meet the weakened
threshold and pass both approximate rows. The upper output for \(S\) exceeds
\(y_S\), and the upper output for \(T\) is at least \(y_S/2>1/3\). This would
give a false certificate of \(v_S<1/3\) with the circuit depth unchanged.

\section{One deterministic game for all the inequalities}
\label{sec:arena}

The offsets in the compiled rows will become edge weights. First we give
every label a common reference. Introduce an origin variable \(o\) and express
each value as a difference from \(h_o\). Before scaling, write \(y_i=h_i-h_o\).
The source box gives the rows
\(h_o\le h_i\) and \(h_i\le h_o+1\). Pin separate terminal variables to
\(h_{\Win}=h_o\) and \(h_{\Lose}=h_o+1\), expressing each equality by both
opposite difference inequalities. The tightened threshold becomes
\(h_o\le h_s-\theta\). Represent each gate label in the same way, as \(h_z-h_o\).
Adding the same real number to every \(h\) preserves the rows, so we can normalize
\(h_o=0\).

Multiply all variables and offsets by \(Q=D\,2^{r+1}\). It clears the threshold
and every gate offset. In the example \(Q=31{,}850{,}496\). The result is a finite
system of weighted maximum inequalities
\begin{equation}\label{eq:rows}
           h_i\le \max_{p\in J_a}\bigl(h_{j(a,p)}+c_{a,p}\bigr)
           \qquad\text{for every row \(a\) with left side \(i\)},
\end{equation}
where \(J_a\) is the nonempty list of terms and \(c_{a,p}\) is an integer.
Here \(h\) denotes the scaled coordinates. A one-term maximum gives a difference
inequality. These are weighted versions of max-atoms. Feasibility over the reals
and the integers now agree. Flooring a real solution coordinatewise preserves
every row, because flooring commutes with a finite maximum and with integer
translation. From an integer solution, subtract \(h_o\) from all coordinates and
divide by \(Q\) to recover a real solution of \(\compiled_r\).

Every row of~\eqref{eq:rows} must hold, and a single term can justify the
inequality in any one row. At a variable, Min picks a row with that left side.
Max sees the row and picks a term.
We use the max-atom game construction of Atserias and Maneva with an additional
start vertex~\cite[\S\S2--4]{am}. In the resulting game \(K\), Min owns a vertex
for every variable. For each row \(a\) with left side \(i\), add a zero-weight
edge from \(i\) to a new vertex \(a\) owned by Max. For each term \(p\in J_a\),
add an edge from \(a\) to \(j(a,p)\) of weight \(c_{a,p}\). Terms with the same
destination may be kept as parallel edges or merged by keeping the largest
offset in that row. Give a variable with no row the tautological row \(h_i\le h_i\).
A fresh Min vertex \(s_0\) has a zero-weight edge to every variable, including the
origin. All vertices have outgoing moves.

For example, consider the two rows at \(A\) in the following table.
\[
\begin{array}{@{}lll@{}}
\toprule
\text{Min's row at \(A\)} & \text{Required inequality} & \text{Max's moves and weights}\\
\midrule
a_1 & h_A\le\max(h_B-1,h_C+1) & B\ (-1),\ C\ (+1)\\
a_2 & h_A\le h_C+2            & C\ (+2)\\
\bottomrule
\end{array}
\]
Take \(h_A=5\), \(h_B=6\), and \(h_C=3\). The same labels satisfy both rows,
and Min's choice on a visit to \(A\) determines which reply works. In \(a_1\),
the term for \(B\) has value \(5\) and the term for \(C\) has value \(4\), so
only \(B\) works. In \(a_2\), the term for \(C\) has value \(5\), so \(C\) works.

\begin{lemma}\label{lem:game}
System~\eqref{eq:rows} has an integer solution if and only if \(\val_K(s_0)\ge0\).
\end{lemma}

\begin{proof}
Given a solution \(h\), choose at each Max row vertex a term whose value is at
least the row's left side. Each two-edge move from variable \(i\) through that row
to \(j\) pays \(c\ge h_i-h_j\). Summing over any number of complete moves
telescopes, giving a lower bound of \(\min_i h_i-\max_i h_i\) on the total weight.
The initial edge and the first edge of any unfinished move have weight zero.
Dividing by the number of edges proves that Max guarantees a nonnegative lower
limiting average.

Conversely, suppose \(\val_K(s_0)\ge0\) and fix an optimal positional Max strategy.
Its selected term at each row leaves a graph in which any negative cycle would
let Min start at a variable on that cycle and repeat it, contradicting the guarantee.
Contract each variable--row--variable move to a single edge \(i\to j\) of its selected weight
\(c\). Add a new sink with a zero-weight edge from every variable to the sink.
Since every variable can reach the sink and the contracted graph has no negative
cycle, the shortest distances \(h_i\) to the sink are finite integers. Following
the edge \(i\to j\) and then a shortest path gives \(h_i\le c+h_j\). Thus the
chosen term satisfies its row, for every row. These distances witness the converse
after the strategy has been fixed.
\end{proof}

\section{Threshold reversal and bit complexity}

Combining Lemma~\ref{lem:reach}, \eqref{eq:tighten}, Proposition~\ref{prop:compile},
and Lemma~\ref{lem:game} gives
\[
                 \val_K(s_0)\ge 0 \quad\Longleftrightarrow\quad v_s<t .
\]
A zero value in \(K\) already certifies the strict inequality, so it must
become negative in \(H\). Put \(N=|V(K)|\). Every negative value that can
occur in \(K\) is at most \(-1/N\), as the cycle argument below shows.
This gap leaves room for the shift by one after reversing the sign and
multiplying by \(N\). Make \(H\) on the same graph with all owners swapped
and with each weight \(w\) replaced by \(-Nw-1\).

Fix a pair of positional strategies \((\sigma,\tau)\) for Max and Min in \(K\).
The play from \(s_0\) eventually repeats a cycle of length at most \(N\), so it has
an ordinary limiting average \(p(\sigma,\tau)\). Using the corresponding pair
\((\tau,\sigma)\) in \(H\) gives the same edge sequence and average
\(-Np(\sigma,\tau)-1\). Positional determinacy in both games~\cite{em,gkk} lets
us calculate the value on these finite strategy sets, giving
\[
 \val_H(s_0)=\max_{\tau}\min_{\sigma}\bigl(-Np(\sigma,\tau)-1\bigr)
            =-N\min_{\tau}\max_{\sigma}p(\sigma,\tau)-1
            =-N\val_K(s_0)-1 .
\]
An optimal positional pair realizes \(\val_K(s_0)\) as a cycle mean. If this value
is negative, its cycle has integer total weight at most \(-1\), so
\(\val_K(s_0)\le -1/N\). Thus \(\val_H(s_0)\ge0\) exactly when
\(\val_K(s_0)<0\). In particular, \(\val_K(s_0)=0\) maps to \(-1\).
Equality \(v_s=t\) rules out the strict certificate, gives \(\val_K(s_0)<0\),
and therefore gives a yes answer in \(H\).

For the size bound, let \(m\) be the number of listed successor occurrences
and \(c\) the number of chance states of \(G\). The integer \(D=b3^{n+1}\)
has \(O(n+\log_2(b+1))\) bits, so \(r\) and \(\bits(Q)=\bits(D)+r+1\) are
also \(O(n+\log_2(b+1))\). Each chance row contributes \(2+2r\) gates of
constant size. The other rows and their terms use \(O(n+m)\) entries, including
the box bounds, and the fresh start adds one edge per variable. Hence \(H\) has
\(O(n+m+cr)\) vertices and edges combined. Before reweighting, all integer
offsets have magnitude at most \(Q\), since \(1/b\le t\le1\) implies
\(0\le\theta\le1\) and every gate shift has magnitude at most one. Each final
weight has magnitude at most \(NQ+1\), and each vertex index uses
\(O(\log_2(N+1))\) bits. The encoder writes the game directly from the
input. The linear programs justify the circuit depth, and the shortest paths in
Lemma~\ref{lem:game} recover labels once a strategy in \(K\) has been fixed.
Computing the parameters by binary arithmetic, writing
the explicit rows and edges, and transforming the weights all take polynomial
time in the full input bit length. This completes the proof of
Theorem~\ref{thm:main}.

\section*{Computational record}

The Python implementation of the reduction, called H4 in the repository, compares
exact rational source decisions with strategy certificates on the emitted graphs.
The replay covered 3,634 game/threshold cases, including the labeled binary games
with one or two nonterminal states and unordered successor pairs. The source answers
agreed with the target certificates and with a separate check of each graph alone
using the Z3 solver. On the equality example the Python implementation,
which retains redundant terminal box bounds, emitted a 193-vertex target and a
verified yes certificate. These are checks of the implementation.
The code, verifier, and report are available at a fixed revision~\cite{artifact}.

At the cited artifact revision \texttt{b72f00b1}, the companion Lean development
checked the strict dual certificate for the supremum of the Bellman iterates,
the bounded rational compiler, the relation between integer
potentials and selected graph cycles, and the sign calculation for positional
cycles with a supplied saddle point. The accompanying written proof supplied
the probabilistic identification of the Bellman limit, the passage from real
strict witnesses to rational certificates, positional determinacy for game
payoffs, and the full encoding and bit-time argument. The Lean compiler used a
general determinant bound and the Python compiler used the game-specific
precision above. The formal correspondence between those implementations
remained open at that revision~\cite{artifact}.

\paragraph{Contributions and use of AI.}
Robert Thomas Collins framed the research question as an exact
threshold-preserving reduction from simple stochastic games to deterministic
mean-payoff games. He directed the investigation and review and worked through
the proof to refine its presentation. OpenAI Codex investigated prior work,
developed the proof, implemented the Python compiler and replay, formalized the
reduction in Lean, and drafted and revised the manuscript. Independent specialist review of the full
reduction remains pending.

\clearpage
\begin{thebibliography}{99}

\bibitem{abgj}
X. Allamigeon, P. Benchimol, S. Gaubert, and M. Joswig.
What tropical geometry tells us about the complexity of linear programming.
\emph{SIAM Review} 63 (2021), 123--164.
\href{https://www.cmap.polytechnique.fr/~gaubert/PAPERS/sigest-paper.pdf\#page=27}{Proposition~20 and equation~(30)}.

\bibitem{am}
A. Atserias and E. Maneva.
Mean-payoff games and the max-atom problem.
Technical report, version of 3 February 2010.
\href{https://www.cs.upc.edu/~atserias/papers/mean-payoff-games-max-atom/mpma.pdf}{Authors' manuscript}.

\bibitem{bnr}
M. Bezem, R. Nieuwenhuis, and E. Rodr\'iguez-Carbonell.
Exponential behaviour of the Butkovi\v{c}--Zimmermann algorithm for solving two-sided
linear systems in max-algebra.
\emph{Discrete Applied Mathematics} 156 (2008), 3506--3509.
\href{https://www.cs.upc.edu/~roberto/papers/dam08.pdf\#page=4}{Construction in Section~3}.

\bibitem{bls}
M. Bodirsky, G. Loho, and M. Skomra.
Reducing stochastic games to semidefinite programming.
In \emph{52nd International Colloquium on Automata, Languages, and Programming},
LIPIcs 334 (2025), article 145.
\href{https://drops.dagstuhl.de/entities/document/10.4230/LIPIcs.ICALP.2025.145}{Conference paper}.

\bibitem{condon}
A. Condon. The complexity of stochastic games.
\emph{Information and Computation} 96 (1992), 203--224.
\href{https://doi.org/10.1016/0890-5401(92)90048-K}{Journal article}.

\bibitem{em}
A. Ehrenfeucht and J. Mycielski. Positional strategies for mean payoff games.
\emph{International Journal of Game Theory} 8 (1979), 109--113.
\href{https://doi.org/10.1007/BF01768705}{Journal article}.

\bibitem{gkk}
V. A. Gurvich, A. V. Karzanov, and L. G. Khachiyan.
Cyclic games and an algorithm to find minimax cycle means in directed graphs.
\emph{USSR Computational Mathematics and Mathematical Physics} 28(5) (1988),
85--91. \href{https://doi.org/10.1016/0041-5553(88)90012-2}{Journal article}.

\bibitem{parys}
P. Parys. Quasipolynomial algorithms for mean-payoff and simple stochastic games.
University of Warsaw, Automata Theory Seminar, 28 October 2026.
\href{https://mimuw.edu.pl/en/seminars/talk/quasipolynomial-algorithms-for-mean-payoff-and-simple-stochastic-games/}{Seminar abstract}.

\bibitem{zp}
U. Zwick and M. Paterson. The complexity of mean payoff games on graphs.
\emph{Theoretical Computer Science} 158 (1996), 343--359.
\href{https://doi.org/10.1016/0304-3975(95)00188-3}{Journal article}.

\bibitem{artifact}
H4 project. Threshold reduction, exact replay, and Lean modules, 2026.
Repository revision \texttt{b72f00b1}.
\href{https://github.com/robbiemu/bayesian-simple-security-games/blob/b72f00b1bb9aebad757f5260a60b3278b2879f0c/notes/H4-ssg-to-mean-payoff.md}{Proof and replay report},
\href{https://github.com/robbiemu/bayesian-simple-security-games/blob/b72f00b1bb9aebad757f5260a60b3278b2879f0c/experiments/h4_ssg_to_mpg_replay.py}{implementation and verifier},
\href{https://github.com/robbiemu/bayesian-simple-security-games/blob/b72f00b1bb9aebad757f5260a60b3278b2879f0c/lean/README.md}{formalization scope}.

\end{thebibliography}
\end{document}
